\section{Der Algorithmus von Kruskal}
Der Algorithmus von Kruskal vom Mathematiker Joseph Bernard Kruskal im Artikel \cite{kruskal_article}.
Der Algorithmus findet für einen ungerichteten endliche Graphen mit Kantenkosten einen minimalen Spannbaum.

\subsection{Was ist ein Spannbaum}
Gegeben ein Graph G(V, E) mit Knotenmenge V und Kantenmenge E, dann ist ein Spannbaum T(V, E') ein Subgraph von G mit E' $\subseteq$ E, sodass: 
\begin{itemize}
    \item alle Knoten aus V, die in G durch einen Kantenzug verbunden sind, auch in T verbunden sind.
    \item in T keine Kreise vorkommen.
\end{itemize}
Wenn also G vollständig verbunden ist, dann ist es auch T und man spricht tatsächlich von einem Spannbaum. 
Besteht G hingegen mehreren Klastern von Konten zwischen denen es keine verbindende Kante gibt, dann gibt es in T entsprechend die selben Klaster und  man  spricht von einem Spannwald.

Ein minimaler Spannbaum ist ein Spannbaum, der die Summe der Kantenkosten minimiert. Das heißt, wenn M(V, E*) ein minimaler Spannbaum von G ist, 
dann muss für jeden anderen Spannbaum T(V, E') von G die Summe der Kantenkosten mindestens so groß sein, wie die von M.
$$\sum_{e \in E^*} c(e) \le \sum_{e \in E'} c(e)$$

\subsection{Der Algorithmus}
Der Algorithmus sortiert zunächst die Kantenliste aufsteigeng nach kosten.
Anschließend iteriert er über die Kanten der Kantenliste, beginnend bei der mit den niedrigsten kosten.
Für jede Kante prüft er nun, ob er die beiden Knoten der Kante schon über einen Kantenzug verbunden hat.
Falls nein übernimmt er die kannte und merkt sich, dass die beiden Knoten jetzt vebunden sind.
Falls ja wird die Kante verworfen.
Zum Schluss gibt er die Liste der übernommenen Kanten zurück.

In Kruskals Worten: 
\begin{quotation}
    Perform the following step as many times as possible: Among the edges of G not yet chosen, choose the shortest edge which does not form any loops with those edges already chosen.\cite{kruskal_article}
\end{quotation}

\subsection{Pseudocode}

\subsection{Waum der Algorithmus funktioniert}
